\appendix
\section*{Appendix B — Operator Algebra}
\addcontentsline{toc}{section}{Appendix B — Operator Algebra}

\section{Operator Framework}

\subsection{Reversible Operators}
Operators are reversible transformations
\[
O : \mathcal{S} \rightarrow \mathcal{S},
\qquad
O^{-1} \text{ exists},
\]
acting on the Fun–Sphere state space $\mathcal{S}$. Reversibility ensures
\[
O^{-1}(O(s)) = s
\]
for all states $s \in \mathcal{S}$.

\subsection{Composition}
Operators compose via
\[
O_j \circ O_i : \mathcal{S} \rightarrow \mathcal{S},
\]
with identity operator $I$ satisfying
\[
I(s) = s.
\]

\subsection{Commutation}
Two operators commute if
\[
O_j \circ O_i = O_i \circ O_j.
\]
Non–commutativity is the source of holonomy curvature.

\section{Algebraic Structures}

\subsection{Operator Group}
The set of reversible operators forms a group:
\begin{itemize}
    \item identity $I$,
    \item inverses $O^{-1}$,
    \item closure under composition,
    \item associativity.
\end{itemize}

\subsection{Operator Path}
A path is a finite sequence
\[
\gamma = (O_1, O_2, \ldots, O_n),
\]
representing a reversible quest.

\subsection{Operator Groupoid}
Paths form a groupoid:
\begin{itemize}
    \item identity path,
    \item inverse path $\gamma^{-1}$,
    \item partial composition,
    \item curvature accumulation.
\end{itemize}

\section{Definitions}

\begin{definition}[Operator Path]
A path $\gamma$ is an ordered sequence of reversible operators acting on
the state space.
\end{definition}

\begin{definition}[Forward Traversal]
Forward traversal applies operators in sequence:
\[
\gamma(s) = O_n(\cdots O_2(O_1(s))\cdots).
\]
\end{definition}

\begin{definition}[Reverse Traversal]
Reverse traversal applies inverses in reverse order:
\[
\gamma^{-1}(s) = O_1^{-1}(O_2^{-1}(\cdots O_n^{-1}(s)\cdots)).
\]
\end{definition}

\begin{definition}[Commutator]
The commutator of two operators is
\[
[O_i, O_j] = O_j \circ O_i - O_i \circ O_j.
\]
\end{definition}

\begin{definition}[Holonomy Contribution]
The holonomy curvature contribution is
\[
\Delta_{\mathrm{hol}}(O_i, O_j) = [O_i, O_j].
\]
\end{definition}

\section{Theorems}

\begin{theorem}[Reversibility]
For any operator path $\gamma$ and state $s$,
\[
\gamma^{-1}(\gamma(s)) = s.
\]
\end{theorem}

\begin{proof}
Each operator $O_i$ has an inverse $O_i^{-1}$ satisfying
$O_i^{-1}(O_i(s)) = s$. Applying inverses in reverse order restores the
original state.
\end{proof}

\begin{theorem}[Curvature from Non–Commutativity]
Curvature is zero if and only if all operators commute.
\end{theorem}

\begin{proof}
Curvature is defined by the holonomy tensor
\[
H_{ij} = [O_i, O_j].
\]
If operators commute, $[O_i, O_j] = 0$ for all pairs, so curvature is
zero. If any pair fails to commute, curvature is non–zero.
\end{proof}

\begin{theorem}[Path Curvature Accumulation]
For a path $\gamma = (O_1, \ldots, O_n)$,
\[
K(\gamma) = \sum_{i < j} [O_i, O_j].
\]
\end{theorem}

\begin{proof}
Curvature accumulates pairwise along the path. The sum over all ordered
pairs yields total curvature.
\end{proof}

\section{Discussion}
This appendix establishes the algebraic foundation for reversible
operators, operator paths, commutation relations, and holonomy curvature.
These structures support the geometric and spectral modules presented in
the main preprint and subsequent appendices.
